\chapter{操作空间动力学 · 阻抗/导纳 · 无源性}
\label{ch:operational-space}

% ════════════════════════════════════════════════════════════════
%  控制部 C3 · 第 O 章：接触交互控制的理论总章。
%  融合 Tutorial 05_运动控制/20_机械臂 F02(数学基础)+F01(导论)+F06(变阻抗/碰撞安全)。
%  记号：右扰动主、se(3)=[rho;phi]、Hamilton 四元数、tilde x = x - x_d、恢复系数 c_r。
%  与 ch:control 严格去重：入门骨架 \cref，本章只做深化/严格化。
% ════════════════════════════════════════════════════════════════

\paragraph{承上：从“控位置”到“控关系”。}
控制部的前几章已把状态空间、能控能观、Lyapunov、综合与最优控制做透（\cref{ch:statespace-basics,ch:controllability-obsv,ch:lyapunov-basics,ch:linear-synthesis,ch:optimal-control-basics}），\cref{ch:control} 又把这些工具落到机器人，给出了关节动力学、计算力矩、PD+重力补偿、操作空间与阻抗控制的\emph{入门骨架}。但那一层把末端当作可以精确定位的刚性执行器；一旦机器人要\emph{接触}世界——插轴、打磨、开门、与人协作——纯位置控制就会“打架”：微小的位置误差被高刚度放大成巨大的接触力。本章正是要把 \cref{ch:control} 只给了结论的那几条线索补成完整的理论链：把多连杆机器人压成一个 6 维等效质量块（Khatib 操作空间动力学），让末端表现为一个可编程的弹簧—阻尼—惯量（Hogan 阻抗），用无源性证明它碰任何被动环境都稳（Colgate--Hogan），再把参数做成时变还保稳（能量罐 / KB 阻尼下界）。这是控制部里“接触交互控制”的理论总章，是卷五《机器人最优控制》中操作空间、阻抗、力控 WBC 与复合系统诸章的共同地基。

\begin{note}
\textbf{本章记号约定（四处与他章不同，先钉死）.}\;
本章是接触李群最深的一章，记号必须与 \cref{ch:lie} 严格对齐，并显式声明与 \cref{ch:control} 的差异：
\begin{enumerate}
  \item \textbf{位姿误差方向}：本章统一取\emph{当前减期望} $\tilde{\mathbf x}=\mathbf x-\mathbf x_d$，恢复力为 $-\mathbf K_d^x\tilde{\mathbf x}$。这与 \cref{eq:ctrl-impedance} 用的 $\mathbf e_x=\mathbf x_d-\mathbf x$ \emph{符号相反}——读那里的公式时按相反约定换号；本章内部一律 current$-$desired。
  \item \textbf{$\se(3)$ 排序}：与 \cref{ch:lie} 一致，$\boldsymbol\xi=[\boldsymbol\rho^\top,\boldsymbol\phi^\top]^\top$（\textbf{平移 $\boldsymbol\rho$ 在前、旋转 $\boldsymbol\phi$ 在后}，$\boldsymbol\rho\neq\mathbf t$）。$6\times6$ 刚度矩阵随之布局：左上块 $\mathbf K_{\rho\rho}$ 管平移、右下块 $\mathbf K_{\phi\phi}$ 管旋转。
  \item \textbf{四元数}：本书用 \textbf{Hamilton} 约定（标量在前 $\mathbf q=[w,\mathbf v]$）；实现库（libfranka/Eigen \texttt{coeffs()}）标量在末，引用其结果时显式换序。
  \item \textbf{恢复系数}：因 $e$ 已被姿态误差 $\mathbf e_R,\mathbf e_p$ 占用，牛顿碰撞恢复系数本章记 $c_r$（取值 $[0,1]$），环境刚度记 $K_e$。
\end{enumerate}
扰动方向全书\emph{右扰动}主，$\mathbf R\,\mathrm{Exp}(\delta\boldsymbol\phi)$、右雅可比 $\mathbf J_r$；引用左扰动 / 异序来源处并列给出经伴随的转换。
\end{note}

\paragraph{启下：八幕主线与历史谱系。}
本章按“动机 $\to$ 操作空间动力学 $\to$ 频域阻抗/导纳 $\to$ 无源性 $\to$ SE(3) 几何 $\to$ 变阻抗保稳 $\to$ 接触/椭球 $\to$ 桥接”八幕展开，每一幕都把 \cref{ch:control} 的入门版严格化。历史上这条线由一串名字串起：Whitney/Draper 实验室的被动柔顺（RCC，1970 年代）$\to$ Salisbury 的主动刚度（1980）$\to$ Mason 的自然/人工约束（1981）$\to$ Hogan 的阻抗哲学（1985）$\to$ Khatib 的操作空间（1987）$\to$ Colgate--Hogan 的耦合稳定（1988）$\to$ DLR/Albu-Sch\"affer 的统一无源框架（2007）$\to$ EPFL/Kronander--Billard 的变阻抗稳定（2016）$\to$ TUM/Haddadin 的力—阻抗统一（2024）。本章把这条 40 年的谱系编织成一条连续的理论叙事。承接本章的“落地”——力位混合、笛卡尔阻抗工程实现、WBC/TSID——见兄弟章 \cref{ch:force-wbc}；复合系统与多模态 MPC 见 \cref{ch:unified-multimodal-mpc}；真臂实测见 \cref{ch:arm-prac}。

% ════════════════════════════════════════════════════════════════
\section{接触交互的根本问题}
\label{sec:os-contact-problem}

\subsection{位置控制为何在接触失败}
\label{sec:os-poscontrol-fail}

\paragraph{动机：一个必然崩溃的打磨实验。}
设想让一台 7 自由度机械臂用砂纸打磨金属工件，按 \cref{ch:control} 的计算力矩律 \cref{eq:ctrl-ctc} 精确跟踪一条贴着工件表面的轨迹。砂纸接触表面的瞬间，会出现接触力瞬间飙升、剧烈振荡、砂纸跳离表面（弹跳）等灾难性现象。根因不在控制律“调得不好”，而在于：\emph{接触时位置误差不再产生恢复运动，而是产生接触力}，位置控制器“不理解力，只知道位置偏了就使劲推”。

\paragraph{接触是一次拓扑突变。}
机械臂从自由空间进入接触，系统发生一次\emph{拓扑突变}：自由度数从 $n$ 降为 $n-k$（$k$ 个约束方向被锁定）。这不是连续渐变，而是瞬态的不连续切换。在接触方向上，把刚性墙环境模型 $\mathbf f_{\rm ext}=K_e(\mathbf x_{\rm wall}-\mathbf x)$ 代入位置控制闭环（\cref{ch:control} 的 \cref{eq:ctrl-manip} 接触项），等效一维动力学为
\begin{equation}
  m_{\rm eff}\ddot x+d_{\rm eff}\dot x+(K_p^{\rm cart}+K_e)x=K_p^{\rm cart}x_d+K_e x_{\rm wall},
  \label{eq:os-contact-stiffness}
\end{equation}
等效刚度变成 $K_p^{\rm cart}+K_e$。对钢铁工件 $K_e$ 可达 $10^6$--$10^8\,$N/m，故自然频率
\begin{equation}
  \omega_n=\sqrt{\frac{K_p^{\rm cart}+K_e}{m_{\rm eff}}}\approx\sqrt{\frac{K_e}{m_{\rm eff}}}
  \label{eq:os-contact-freq}
\end{equation}
远超控制器带宽，导致不可控的高频振荡；而稳态接触力 $f_{\rm contact}\approx K_e(x_{\rm wall}-x_d)$ \emph{完全由环境刚度与位置误差决定}——在 $K_e=10^6\,$N/m 下，$0.1\,$mm 的标定误差就产生约 $100\,$N 的力。接触瞬间的冲击力按量级估计为
\begin{equation}
  f_{\rm impact}\sim\frac{m_{\rm eff}\,v_0}{\Delta t},
  \label{eq:os-impact}
\end{equation}
$\Delta t\sim1$--$5\,$ms 为接触持续时间。以 $m_{\rm eff}\approx3\,$kg、$v_0=0.1\,$m/s、$\Delta t\approx2\,$ms 估得峰值约 $150\,$N。须强调：该数值\emph{强依赖}有效质量、恢复系数与材料，只能作量级参考，不可当作可引用的结论。

\begin{pitfall}[概念误区：用高增益位置控制器去“顶”接触力]
新手常想“接触力太大就把位置增益调软一点”，却仍用纯位置环。问题是纯位置控制的等效刚度由反馈增益隐式决定且常很高，接触时位置误差 $\to$ 巨大力；调低增益又牺牲自由空间精度。根本原因：位置与力在接触方向上\emph{不能同时独立指定}——环境约束把二者耦合。\cref{ch:control} 的 \cref{eq:ctrl-impedance} 已给出正解（显式设定力—运动关系）；本章把这条线索补成完整的操作空间动力学与无源性理论。
\end{pitfall}

\subsection{Mason 的自然约束与人工约束}
\label{sec:os-mason}

\paragraph{接触把 6 维操作空间正交分解。}
在设计任何力控律之前，需要 Mason（1981）的约束空间分析\cite{mason1981compliance} 回答：哪些自由度该控位置、哪些该控力？Mason 的核心观察是，末端接触环境时，6 维操作空间被分解为两个正交子空间：
\begin{itemize}
  \item \textbf{自然约束}（natural constraints）：由接触几何\emph{决定}、不可违反。沿法向不能运动（位移被锁），但可施力。
  \item \textbf{人工约束}（artificial constraints）：由工程师\emph{选择}、用控制实现。沿切向位置需跟踪，而切向力设为零或受控。
\end{itemize}
形式化地，引入对角的\emph{选择矩阵} $\mathbf S$（$0/1$ 对角，第 $i$ 个对角元为 $1$ 表示该方向控力），其补 $\bar{\mathbf S}=\mathbf I-\mathbf S$ 选位置方向，满足
\begin{equation}
  \mathbf S+\bar{\mathbf S}=\mathbf I,\qquad \mathbf S\bar{\mathbf S}=\mathbf 0.
  \label{eq:os-selection}
\end{equation}
力控方向用 $\mathbf S$ 选出、位置控制方向用 $\bar{\mathbf S}$ 选出，二者在一个控制律里互补——这正是 Raibert--Craig（1981）混合力位控制\cite{raibert1981hybrid} 的出发点。本章只点出选择矩阵这一几何地基，\emph{完整}的混合力位（含并行力位、Chiaverini 结构）留给 \cref{ch:force-wbc}。

\subsection{Hogan 哲学：控关系，不控变量}
\label{sec:os-hogan}

\cref{ch:control} 的 \cref{def:ctrl-impedance} 已给出机械阻抗（运动 $\to$ 力的算子）与导纳（其逆）的定义，及 Hogan 的核心论断“两个互连子系统不能都是导纳，故机械手应表现为阻抗、环境为导纳”。这里只补一条\emph{能量端口}视角的论证，作为后文无源性的引子：把机器人与环境的交互看作通过一个端口交换能量，端口功率为
\begin{equation}
  P=\mathbf f^\top\mathbf v,
  \label{eq:os-port-power}
\end{equation}
$\mathbf f$ 是势变量（effort）、$\mathbf v$ 是流变量（flow）。在一个端口上，因果关系必须相容：若两端都企图“给定运动、由对方决定力”（都是导纳），接触瞬间速度无法相容（刚性碰撞速度跳变）；若两端都“给定力”，则位置漂移。唯一相容的配置是一端为阻抗、一端为导纳。Hogan 据此主张机械手呈现阻抗——这一哲学贯穿全章，下文将反复以 \cref{eq:os-port-power} 的功率作为能量记账的端口变量。

% ════════════════════════════════════════════════════════════════
\section{操作空间动力学}
\label{sec:os-opspace-dynamics}

\paragraph{动机：把多连杆压成一个 6D 质量块。}
力控律要在\emph{笛卡尔空间}（任务在那里定义，如“沿桌面以 $20\,$N 法向力滑动”）设计，但动力学天然写在关节空间。\cref{ch:control} 的 \cref{eq:ctrl-osc} 已给出 Khatib 操作空间动力学的\emph{结论式} $\boldsymbol\Lambda\ddot{\mathbf x}+\boldsymbol\mu+\mathbf p=\mathbf F$ 与 $\boldsymbol\Lambda=(\mathbf J\mathbf M^{-1}\mathbf J^\top)^{-1}$，却把推导跳过了。本节补全这条推导，并揭示它为何让阻抗控制“自然”工作：操作空间动力学本质上把整条多连杆机械臂\emph{压缩}成一个等效的 6D 质量块——你从外部推末端，末端的加速度由“等效质量”$\boldsymbol\Lambda$ 决定，但 $\boldsymbol\Lambda$ 随构型变化（某些姿态“感觉重”、某些“感觉轻”）；$\boldsymbol\mu$ 是“等效科氏力”，$\mathbf p$ 是“等效重力”。阻抗控制就是把这个真实质量块替换成一个\emph{设计的}虚拟弹簧—阻尼—惯量。

\subsection{Khatib 投影：完整推导}
\label{sec:os-khatib-derivation}

\begin{theorem}[操作空间动力学（Khatib 1987）]
\label{thm:os-opspace}
设 $n$ 自由度机械臂的关节动力学为 $\mathbf M(\mathbf q)\ddot{\mathbf q}+\mathbf C(\mathbf q,\dot{\mathbf q})\dot{\mathbf q}+\mathbf g(\mathbf q)=\boldsymbol\tau+\mathbf J^\top\mathbf f_{\rm ext}$，任务雅可比 $\dot{\mathbf x}=\mathbf J\dot{\mathbf q}$（$\mathbf J\in\mathbb R^{m\times n}$，$m\le n$，满行秩）。则末端任务坐标满足 \cref{eq:ctrl-osc} 的操作空间动力学，其中
\begin{equation}
  \boldsymbol\Lambda=(\mathbf J\mathbf M^{-1}\mathbf J^\top)^{-1},\quad
  \bar{\mathbf J}=\mathbf M^{-1}\mathbf J^\top\boldsymbol\Lambda,\quad
  \boldsymbol\mu=\bar{\mathbf J}^\top\mathbf C\dot{\mathbf q}-\boldsymbol\Lambda\dot{\mathbf J}\dot{\mathbf q},\quad
  \mathbf p=\bar{\mathbf J}^\top\mathbf g,\quad
  \mathbf F=\bar{\mathbf J}^\top\boldsymbol\tau+\mathbf f_{\rm ext}.
  \label{eq:os-opspace-terms}
\end{equation}
\end{theorem}
\begin{proof}
分六步。\textbf{Step 1}（运动学）：对 $\dot{\mathbf x}=\mathbf J\dot{\mathbf q}$ 求导得 $\ddot{\mathbf x}=\mathbf J\ddot{\mathbf q}+\dot{\mathbf J}\dot{\mathbf q}$。
\textbf{Step 2}（加速度分解）：对冗余臂（$n>m$），关节加速度分解为任务分量与零空间分量
\begin{equation}
  \ddot{\mathbf q}=\bar{\mathbf J}(\ddot{\mathbf x}-\dot{\mathbf J}\dot{\mathbf q})+(\mathbf I-\bar{\mathbf J}\mathbf J)\ddot{\mathbf q}_{\rm null},
  \label{eq:os-accel-decomp}
\end{equation}
其中 $\bar{\mathbf J}$ 是待定的广义逆。\textbf{Step 3}（定 $\bar{\mathbf J}$）：把 \cref{eq:os-accel-decomp} 代入关节动力学并左乘 $\bar{\mathbf J}^\top$，零空间项为 $\bar{\mathbf J}^\top\mathbf M(\mathbf I-\bar{\mathbf J}\mathbf J)\ddot{\mathbf q}_{\rm null}$。要使零空间与任务空间\emph{动力学解耦}，须令该项对任意 $\ddot{\mathbf q}_{\rm null}$ 为零，即 $\bar{\mathbf J}^\top\mathbf M(\mathbf I-\bar{\mathbf J}\mathbf J)=\mathbf 0$。代入 $\boldsymbol\Lambda=(\mathbf J\mathbf M^{-1}\mathbf J^\top)^{-1}$ 验证：唯一满足该约束的广义逆是\emph{动力学一致伪逆}
\begin{equation}
  \bar{\mathbf J}=\mathbf M^{-1}\mathbf J^\top(\mathbf J\mathbf M^{-1}\mathbf J^\top)^{-1}=\mathbf M^{-1}\mathbf J^\top\boldsymbol\Lambda.
  \label{eq:os-jbar}
\end{equation}
\textbf{Step 4}（验证 $\bar{\mathbf J}^\top\mathbf M\bar{\mathbf J}=\boldsymbol\Lambda$）：
\[
  \bar{\mathbf J}^\top\mathbf M\bar{\mathbf J}=\boldsymbol\Lambda\mathbf J\mathbf M^{-1}\,\mathbf M\,\mathbf M^{-1}\mathbf J^\top\boldsymbol\Lambda=\boldsymbol\Lambda(\mathbf J\mathbf M^{-1}\mathbf J^\top)\boldsymbol\Lambda=\boldsymbol\Lambda\boldsymbol\Lambda^{-1}\boldsymbol\Lambda=\boldsymbol\Lambda.
\]
\textbf{Step 5}（验证 $\mathbf J\bar{\mathbf J}=\mathbf I$）：$\mathbf J\bar{\mathbf J}=\mathbf J\mathbf M^{-1}\mathbf J^\top(\mathbf J\mathbf M^{-1}\mathbf J^\top)^{-1}=\mathbf I$，故 $\bar{\mathbf J}^\top\mathbf J^\top=(\mathbf J\bar{\mathbf J})^\top=\mathbf I$。
\textbf{Step 6}（整理）：把上述代入左乘 $\bar{\mathbf J}^\top$ 后的方程，第一项 $\bar{\mathbf J}^\top\mathbf M\bar{\mathbf J}(\ddot{\mathbf x}-\dot{\mathbf J}\dot{\mathbf q})=\boldsymbol\Lambda\ddot{\mathbf x}-\boldsymbol\Lambda\dot{\mathbf J}\dot{\mathbf q}$，零空间项依设计为零，外力项 $\bar{\mathbf J}^\top\mathbf J^\top\mathbf f_{\rm ext}=\mathbf f_{\rm ext}$，移项即得 \cref{eq:os-opspace-terms} 的各项。
\end{proof}

\begin{insight}
关节力矩由对偶映射 $\boldsymbol\tau=\mathbf J^\top\mathbf F$ 给出（虚功原理 $\boldsymbol\tau^\top\delta\mathbf q=\mathbf F^\top\delta\mathbf x$ 与 $\delta\mathbf x=\mathbf J\delta\mathbf q$ 立得）。任务控制力取 $\mathbf F=\hat{\boldsymbol\Lambda}\mathbf F^*+\hat{\boldsymbol\mu}+\hat{\mathbf p}$（$\hat{\cdot}$ 为模型估计）即抵消非线性，给出解耦行为 $\ddot{\mathbf x}=\mathbf F^*$——这是 \cref{eq:ctrl-osc} 末解耦律的来历，也是后文阻抗律的载体。
\end{insight}

\subsection{科氏项 \texorpdfstring{$\boldsymbol\mu$}{mu} 的物理拆解}
\label{sec:os-mu}

$\boldsymbol\mu=\bar{\mathbf J}^\top\mathbf C\dot{\mathbf q}-\boldsymbol\Lambda\dot{\mathbf J}\dot{\mathbf q}$ 是教材最常“跳过”的一项，两部分各有清晰来源：$\bar{\mathbf J}^\top\mathbf C\dot{\mathbf q}$ 是关节空间科氏/离心力在操作空间的投影；$-\boldsymbol\Lambda\dot{\mathbf J}\dot{\mathbf q}$ 是雅可比时变导致的虚拟力。后者常被误解为“关节加速度的贡献”，实则是\emph{参考系变化率}的体现：即使关节速度恒定（$\ddot{\mathbf q}=\mathbf 0$），末端加速度 $\ddot{\mathbf x}=\dot{\mathbf J}\dot{\mathbf q}\neq\mathbf 0$，因为 $\mathbf J$ 随构型变化。这与旋转参考系中的科氏力同源——来源于参考系本身在变。高速运动时若漏掉 $\dot{\mathbf J}\dot{\mathbf q}$ 补偿，末端会偏离期望轨迹。

\subsection{动力学一致伪逆 vs Moore--Penrose 伪逆}
\label{sec:os-jbar-vs-pinv}

\begin{table}[H]\centering\small
\caption{动力学一致伪逆 $\bar{\mathbf J}$ 与 Moore--Penrose 伪逆 $\mathbf J^+$ 的对照}
\label{tab:os-jbar-vs-pinv}
\begin{tabular}{p{0.20\textwidth} p{0.36\textwidth} p{0.36\textwidth}}
\toprule
性质 & 动力学一致 $\bar{\mathbf J}=\mathbf M^{-1}\mathbf J^\top\boldsymbol\Lambda$ & Moore--Penrose $\mathbf J^+=\mathbf J^\top(\mathbf J\mathbf J^\top)^{-1}$ \\
\midrule
加权范数 & 最小动能 $\dot{\mathbf q}^\top\mathbf M\dot{\mathbf q}$ & 最小欧氏 $\|\dot{\mathbf q}\|_2^2$ \\
零空间解耦 & \textbf{是}：零空间力矩不影响 $\ddot{\mathbf x}$ & 否：经惯量耦合“泄漏”到 $\ddot{\mathbf x}$ \\
适用层级 & 力矩级控制（阻抗 / WBC） & 速度级控制（逆运动学） \\
\bottomrule
\end{tabular}
\end{table}

\begin{insight}
$\bar{\mathbf J}$ 与 $\mathbf J^+$ 的差异\emph{不是}精度问题而是\emph{物理正确性}问题。在阻抗控制中若用 $\mathbf J^+$ 做零空间投影，零空间力矩会经惯量耦合泄漏到末端——这正是许多“阻抗控制末端抖动”的数学根因（真臂实测见 \cref{sec:ap-illcond}）。只有 $\bar{\mathbf J}$ 保证零空间与任务空间在\emph{动能度量}下解耦。
\end{insight}

\subsection{零空间投影：幂等性与动力学一致性}
\label{sec:os-nullspace}

\begin{theorem}[零空间投影矩阵的性质]
\label{thm:os-nullspace}
设 $\mathbf N=\mathbf I-\bar{\mathbf J}\mathbf J$。则 (i) 幂等 $\mathbf N^2=\mathbf N$；(ii) 任务零化 $\mathbf J\mathbf N=\mathbf 0$；(iii) 动力学一致 $\mathbf N^\top\mathbf M=\mathbf M\mathbf N$。
\end{theorem}
\begin{proof}
(i) 用 $\mathbf J\bar{\mathbf J}=\mathbf I$：
\[
  \mathbf N^2=(\mathbf I-\bar{\mathbf J}\mathbf J)(\mathbf I-\bar{\mathbf J}\mathbf J)=\mathbf I-2\bar{\mathbf J}\mathbf J+\bar{\mathbf J}(\mathbf J\bar{\mathbf J})\mathbf J=\mathbf I-2\bar{\mathbf J}\mathbf J+\bar{\mathbf J}\mathbf J=\mathbf N.
\]
(ii) $\mathbf J\mathbf N=\mathbf J-(\mathbf J\bar{\mathbf J})\mathbf J=\mathbf J-\mathbf J=\mathbf 0$。
(iii) 一侧 $\mathbf N^\top\mathbf M=(\mathbf I-\mathbf J^\top\bar{\mathbf J}^\top)\mathbf M=\mathbf M-\mathbf J^\top\boldsymbol\Lambda\mathbf J\mathbf M^{-1}\mathbf M=\mathbf M-\mathbf J^\top\boldsymbol\Lambda\mathbf J$（用 $\bar{\mathbf J}^\top=\boldsymbol\Lambda\mathbf J\mathbf M^{-1}$）；另一侧 $\mathbf M\mathbf N=\mathbf M-\mathbf M\bar{\mathbf J}\mathbf J=\mathbf M-\mathbf M\mathbf M^{-1}\mathbf J^\top\boldsymbol\Lambda\mathbf J=\mathbf M-\mathbf J^\top\boldsymbol\Lambda\mathbf J$，二者相等。
\end{proof}

性质 (iii) 说 $\mathbf N$ 是 $\mathbf M$-正交投影：它在\emph{动能度量}下把运动分解为任务与零空间分量，两者在能量意义上完全解耦。若改用 $\mathbf J^+$，得到的投影只是欧氏正交、动能度量下不解耦，零空间力矩遂经惯量泄漏到任务空间——这与 \cref{sec:os-jbar-vs-pinv} 同一根因。在此基础上，多优先级控制律（Sentis--Khatib 2005）层叠为
\begin{equation}
  \boldsymbol\tau=\mathbf J_1^\top\mathbf F_1+\mathbf N_1^\top\mathbf J_2^\top\mathbf F_2+\mathbf N_2^\top\boldsymbol\tau_{\rm null},
  \label{eq:os-multipriority}
\end{equation}
$\mathbf F_1$ 是最高优先级任务力（如末端阻抗），$\mathbf F_2$ 是次级任务（关节中心化、操作度优化等），逐级在更高级的零空间内施加。这里只点出零空间投影的数学；\emph{完整}的多优先级 WBC 分层（HQP/TSID）留给 \cref{ch:force-wbc}，足式机器人的同一数学见 \cref{ch:quad_wbc}。

\subsection{奇异性与阻尼最小二乘}
\label{sec:os-dls}

在奇异位形附近 $\boldsymbol\Lambda=(\mathbf J\mathbf M^{-1}\mathbf J^\top)^{-1}$ 的条件数趋于无穷，阻抗律 $\boldsymbol\tau=\mathbf J^\top\boldsymbol\Lambda(\cdots)$ 会命令无穷大力矩。阻尼最小二乘（DLS）正则化
\begin{equation}
  \boldsymbol\Lambda_{\rm DLS}=(\mathbf J\mathbf M^{-1}\mathbf J^\top+\lambda^2\mathbf I)^{-1}
  \label{eq:os-dls}
\end{equation}
以牺牲退化方向的末端精度为代价，保证力矩有限（DLS 的最小二乘背景见 \cref{ch:nlopt}）。可变阻尼策略令 $\lambda^2$ 在最小奇异值 $\sigma_{\min}(\mathbf J)$ 大于阈值 $\sigma_0$ 时为零、低于阈值时随 $\bigl(1-(\sigma_{\min}/\sigma_0)^2\bigr)$ 平滑增大。奇异点附近力椭球退化（某些方向施力能力趋零），阻抗控制无法在退化方向产生足够恢复力，故力控轨迹规划应预先保证 $\sigma_{\min}$ 不过小（见 \cref{sec:os-ellipsoid}）。

% ════════════════════════════════════════════════════════════════
\section{阻抗与导纳的频域理论}
\label{sec:os-frequency}

\paragraph{动机：把“弹簧旋钮”放到频率轴上看。}
\cref{ch:control} 的 \cref{eq:ctrl-impedance} 已给出目标阻抗 $\mathbf M_d\ddot{\tilde{\mathbf x}}+\mathbf D_d\dot{\tilde{\mathbf x}}+\mathbf K_d^x\tilde{\mathbf x}=\mathbf f_{\rm ext}$ 与“刚度是柔顺旋钮”的直觉。本节把它搬到频域，揭示阻抗在不同频段的行为、二端口网络的统一语言，以及环境耦合后的闭环特性——这是设计力控带宽与判断稳定性的工具。

\subsection{刚度/阻尼/惯量矩阵的结构}
\label{sec:os-matrices}

目标阻抗的三个矩阵都有结构约束：$\mathbf K_d^x,\mathbf D_d,\mathbf M_d$ 须对称正定（$\mathbf K_d^x$ 对称是 $\frac12\tilde{\mathbf x}^\top\mathbf K_d^x\tilde{\mathbf x}$ 成为合法势能的前提）。Albu-Sch\"affer 等（2007）的\emph{双对角化}\cite{albuschaffer2007unified} 给出一种工程上优雅的解耦：同时合同对角化 $\mathbf M_d$ 与 $\mathbf K_d^x$，把 $6$ 个耦合方向解耦成 $6$ 个独立的二阶环，再逐环按临界阻尼配 $\mathbf D_d$。临界阻尼矩阵（$\zeta=1$，使各模态无超调）取
\begin{equation}
  \mathbf D_d=2\,\mathbf M_d^{1/2}\bigl(\mathbf M_d^{-1/2}\mathbf K_d^x\mathbf M_d^{-1/2}\bigr)^{1/2}\mathbf M_d^{1/2},
  \label{eq:os-critical-damping}
\end{equation}
对标量退化为熟悉的 $d=2\sqrt{km}$。期望惯量 $\mathbf M_d$ 有三种选择策略：取 $\mathbf M_d=\boldsymbol\Lambda$（不整形惯性，\emph{无需力传感}，律退化为任务空间 PD）、取恒定对角阵（整形惯性、需力反馈）、或取零（纯弹簧—阻尼）。

\subsection{阻抗、导纳与柔顺度的传递函数}
\label{sec:os-transfer}

单自由度目标阻抗的传递函数（以速度为流变量、力为势变量）为
\begin{equation}
  Z(s)=\frac{F(s)}{\dot{\tilde X}(s)}=M_d s+D_d+\frac{K_d}{s},\qquad
  Y(s)=Z(s)^{-1},\qquad
  G(s)=\frac{\tilde X(s)}{F(s)}=\frac{1}{M_d s^2+D_d s+K_d},
  \label{eq:os-Zs}
\end{equation}
$Y(s)$ 是导纳、$G(s)$ 是柔顺度（位移柔度）。其 Bode 幅相有三段清晰的渐近行为：低频弹簧主导（$Z\approx K_d/s$，相位 $-90^\circ$）、共振频率 $\omega_n=\sqrt{K_d/M_d}$ 附近阻尼主导（$Z\approx D_d$，相位 $0^\circ$）、高频惯量主导（$Z\approx M_d s$，相位 $+90^\circ$）。多自由度时各轴带宽不必相同——\cref{eq:os-critical-damping} 的对角化正是为了让每个模态独立设带宽。SISO 二阶系统的指标与裕度细节（上升时间、超调、相位裕度）见 \cref{ch:control} 的 \cref{eq:ctrl-2nd-order-spec,def:ctrl-margins}，此处不重述。

\begin{figure}[ht]
\centering
\begin{tikzpicture}[>=Stealth, font=\small]
  % axes
  \draw[->,thick] (0,0) -- (8.4,0) node[right]{$\lg\omega$};
  \draw[->,thick] (0,0) -- (0,3.4) node[above]{$|Z(j\omega)|$ (dB)};
  % spring asymptote (slope -20 dB/dec)
  \draw[third,very thick] (0.3,3.0) -- (3.2,0.9);
  \node[third] at (1.3,2.55) {$K_d/s$ ($-90^\circ$)};
  % damping plateau
  \draw[main,very thick] (3.2,0.9) -- (5.2,0.9);
  \node[main] at (4.2,1.25) {$D_d$ ($0^\circ$)};
  % inertia asymptote (slope +20 dB/dec)
  \draw[second,very thick] (5.2,0.9) -- (8.1,3.0);
  \node[second] at (7.0,2.55) {$M_d s$ ($+90^\circ$)};
  % resonance mark
  \draw[gray,dashed] (4.2,0) -- (4.2,0.9);
  \node[gray,below] at (4.2,0) {$\omega_n=\sqrt{K_d/M_d}$};
\end{tikzpicture}
\caption{阻抗 $Z(j\omega)$ 的 Bode 幅值渐近：低频弹簧（$-20\,$dB/dec）、共振平台（阻尼主导）、高频惯量（$+20\,$dB/dec）。相位从 $-90^\circ$ 经 $0^\circ$ 升到 $+90^\circ$。}
\label{fig:os-bode}
\end{figure}

\subsection{二端口网络与统一框架}
\label{sec:os-twoport}

把“机器人 + 控制器”抽象为电气网络理论里的\emph{二端口网络}：端口 1 接控制器/操作者、端口 2 接环境，每个端口有一对力—速度变量。力—速度与电压—电流对偶（力对应势变量、速度对应流变量、$P=\mathbf f^\top\mathbf v$ 对应 $VI$）。阻抗矩阵、导纳矩阵与混合矩阵分别描述
\begin{equation}
  \begin{bmatrix}\mathbf F_1\\\mathbf F_2\end{bmatrix}=\mathbf Z\begin{bmatrix}\mathbf V_1\\\mathbf V_2\end{bmatrix},\qquad
  \begin{bmatrix}\mathbf V_1\\\mathbf V_2\end{bmatrix}=\mathbf Y\begin{bmatrix}\mathbf F_1\\\mathbf F_2\end{bmatrix}=\mathbf Z^{-1}\begin{bmatrix}\mathbf F_1\\\mathbf F_2\end{bmatrix},\qquad
  \begin{bmatrix}\mathbf F_1\\\mathbf V_2\end{bmatrix}=\mathbf H\begin{bmatrix}\mathbf V_1\\\mathbf F_2\end{bmatrix}.
  \label{eq:os-twoport}
\end{equation}
混合矩阵 $\mathbf H$ 在遥操作里最常用：理想\emph{透明度}要求 $\mathbf H=\bigl[\begin{smallmatrix}\mathbf 0&\mathbf I\\\mathbf I&\mathbf 0\end{smallmatrix}\bigr]$，即操作者直接“感受到”环境（遥操作二端口的完整理论属机械臂部双臂/遥操作专题，本章只指针）。这套语言把所有力控范式统一为“端口 2 行为的选择”：位置控制是 $K_d\to\infty$（无穷大输出阻抗、刚性速度源）、纯力控是 $K_d,D_d,M_d\to0$（零阻抗、纯力源），阻抗控制是介于二者之间、$Z_{22}$ 可编程的\emph{连续谱}。这正是 Hogan 所谓“统一框架”的精确含义。

\subsection{环境耦合后的闭环}
\label{sec:os-coupled}

把目标阻抗 $G(s)$ 与刚度为 $K_e$ 的弹性环境耦合（环境施力 $f_{\rm ext}=-K_e\tilde x$，单自由度），闭环特征多项式为 $M_d s^2+D_d s+(K_d+K_e)$，故闭环自然频率与阻尼比
\begin{equation}
  \omega_{\rm cl}=\sqrt{\frac{K_d+K_e}{M_d}},\qquad
  \zeta_{\rm cl}=\frac{D_d}{2\sqrt{(K_d+K_e)M_d}}.
  \label{eq:os-coupled-freq}
\end{equation}
注意：环境刚度 $K_e$ \emph{并联}进总刚度，故耦合后频率升高、阻尼比下降——这解释了为何接触刚性环境时易振荡（呼应 \cref{eq:os-contact-freq}），也提示变阻抗时须同步加阻尼（见 \cref{sec:os-variable}）。给定期望偏置 $x_d$，稳态接触力（令 $s\to0$）为
\begin{equation}
  f_{ss}=\frac{K_e K_d}{K_d+K_e}\,x_d,
  \label{eq:os-fss}
\end{equation}
其结构与两弹簧串联的等效刚度一致。一条重要设计原则随之而来：当 $K_d\ll K_e$（控制器“软”），$f_{ss}\approx K_d x_d$，\emph{接触力主要由控制器刚度决定}、几乎与未知环境刚度无关——这正是用低刚度阻抗实现可控接触力的根据。

\subsection{阻抗 vs 导纳的因果对偶与 UFIC}
\label{sec:os-causality}

\begin{table}[H]\centering\small
\caption{阻抗控制与导纳控制的因果对偶}
\label{tab:os-imp-adm}
\begin{tabular}{p{0.18\textwidth} p{0.37\textwidth} p{0.37\textwidth}}
\toprule
维度 & 阻抗控制 & 导纳控制 \\
\midrule
端口行为 & 测运动、输出力（力矩） & 测力、输出运动（位置参考） \\
硬件接口 & 力矩源（无需力传感的特例存在） & 力传感器 + 内层位置环 \\
适合机器人 & 轻质、反驱（背隙小） & 刚性、高减速比 \\
带宽 & 高 & 受内层位置环限制 \\
\bottomrule
\end{tabular}
\end{table}

阻抗与导纳并非互斥的两类，而是同一连续谱的两端。Haddadin 与 Shahriari（2024）的\emph{统一力—阻抗控制}（UFIC）\cite{shahriari2024unified} 用一个标量 $\alpha\in[0,1]$ 在“纯力跟踪”与“纯阻抗”之间连续插值，并配能量罐保证对任意被动环境的稳定（其能量罐机制见 \cref{sec:os-energytank}）。本章把它作为“力—阻抗连续谱”的前沿一段；学习型力控（RL 输出变阻抗）的正式理论留给 \cref{ch:force-wbc} 与强化学习运控部。

% ════════════════════════════════════════════════════════════════
\section{无源性理论}
\label{sec:os-passivity}

\paragraph{动机：环境模型未知，仍要保证稳定。}
接触刚度 $K_e$ 跨 $10^2$（泡沫）到 $10^8$（钢铁）六个数量级。若稳定性依赖已知 $K_e$，控制器就得为每种环境重设计——不可行。无源性提供\emph{无模型}的稳定保证：只要控制器无源，它与\emph{任何}被动环境的耦合都稳定。本节是本章的理论高峰。一般的无源性 / Port-Hamiltonian / LaSalle 系统理论见 \cref{ch:lyapunov-stability}；这里只取接触交互控制所需，并把正实/KYP 引理在阻抗里的角色\emph{证一次}，鲁棒控制部的 IQC 视角（\cref{ch:robust-hinf}）引用而不重证。

\subsection{无源性的严格定义}
\label{sec:os-passivity-def}

\begin{definition}[输入—输出无源与严格输出无源]
\label{def:os-passivity}
系统 $\Sigma$ 以 $\mathbf u(t)$ 为输入、$\mathbf y(t)$ 为输出。若存在常数 $\beta\ge0$ 使
\begin{equation}
  \int_0^T\mathbf u^\top\mathbf y\,\mathrm dt\ge-\beta,\qquad\forall T\ge0,
  \label{eq:os-passive-integral}
\end{equation}
则称 $\Sigma$ \emph{无源}。等价地，若存在存储函数 $V(\mathbf x)\ge0$ 满足\emph{耗散不等式} $\dot V\le\mathbf u^\top\mathbf y$，则无源。若更强地 $\dot V\le\mathbf u^\top\mathbf y-\epsilon\|\mathbf y\|^2$（$\epsilon>0$），则称\emph{严格输出无源}。
\end{definition}

物理含义：端口功率 $\mathbf u^\top\mathbf y$ 的累计能量恒大于某有限负值——系统至多释放有限的初始储能，不能凭空产生无限能量。阻抗控制器的阻尼项 $\mathbf D_d$ 正是提供严格输出无源所需的“主动耗散”。

\begin{insight}
电路类比把无源性讲透：纯电阻—电感—电容网络（无源）可安全接到任何负载、不自激；含放大器（有源）则可能对某些负载自激。无源控制器是“安全的接口”。在阻抗里，$K_d$ 是电容（存弹性势能）、$D_d$ 是电阻（耗能）、$M_d$ 是电感（存动能）。
\end{insight}

\subsection{阻抗控制无源性的完整证明}
\label{sec:os-impedance-passivity}

\begin{theorem}[标准阻抗控制器严格输出无源]
\label{thm:os-impedance-passive}
令末端行为满足 $\mathbf M_d\ddot{\tilde{\mathbf x}}+\mathbf D_d\dot{\tilde{\mathbf x}}+\mathbf K_d^x\tilde{\mathbf x}=\mathbf f_{\rm ext}$，端口取 $\mathbf u=\mathbf f_{\rm ext}$、$\mathbf y=\dot{\tilde{\mathbf x}}$。则该系统严格输出无源，$\epsilon=\lambda_{\min}(\mathbf D_d)$。
\end{theorem}
\begin{proof}
分五步。\textbf{Step 1}（存储函数）：取
\begin{equation}
  V=\tfrac12\dot{\tilde{\mathbf x}}^\top\mathbf M_d\dot{\tilde{\mathbf x}}+\tfrac12\tilde{\mathbf x}^\top\mathbf K_d^x\tilde{\mathbf x},
  \label{eq:os-storage}
\end{equation}
$\mathbf M_d,\mathbf K_d^x\succ0\Rightarrow V\ge0$，且 $V=0\iff\tilde{\mathbf x}=\mathbf 0,\dot{\tilde{\mathbf x}}=\mathbf 0$。
\textbf{Step 2}（求导）：$\dot V=\dot{\tilde{\mathbf x}}^\top\mathbf M_d\ddot{\tilde{\mathbf x}}+\tilde{\mathbf x}^\top\mathbf K_d^x\dot{\tilde{\mathbf x}}$。
\textbf{Step 3}（代入动力学 $\mathbf M_d\ddot{\tilde{\mathbf x}}=\mathbf f_{\rm ext}-\mathbf D_d\dot{\tilde{\mathbf x}}-\mathbf K_d^x\tilde{\mathbf x}$）：
\[
  \dot V=\dot{\tilde{\mathbf x}}^\top(\mathbf f_{\rm ext}-\mathbf D_d\dot{\tilde{\mathbf x}}-\mathbf K_d^x\tilde{\mathbf x})+\tilde{\mathbf x}^\top\mathbf K_d^x\dot{\tilde{\mathbf x}}.
\]
\textbf{Step 4}（交叉项相消）：因 $\mathbf K_d^x$ 对称，标量 $\dot{\tilde{\mathbf x}}^\top\mathbf K_d^x\tilde{\mathbf x}=\tilde{\mathbf x}^\top\mathbf K_d^x\dot{\tilde{\mathbf x}}$，故 $-\dot{\tilde{\mathbf x}}^\top\mathbf K_d^x\tilde{\mathbf x}+\tilde{\mathbf x}^\top\mathbf K_d^x\dot{\tilde{\mathbf x}}=0$，余
\[
  \dot V=\dot{\tilde{\mathbf x}}^\top\mathbf f_{\rm ext}-\dot{\tilde{\mathbf x}}^\top\mathbf D_d\dot{\tilde{\mathbf x}}=\mathbf u^\top\mathbf y-\mathbf y^\top\mathbf D_d\mathbf y.
\]
\textbf{Step 5}（结论）：$\mathbf D_d\succ0\Rightarrow\mathbf y^\top\mathbf D_d\mathbf y\ge\lambda_{\min}(\mathbf D_d)\|\mathbf y\|^2$，故
\begin{equation}
  \dot V\le\mathbf u^\top\mathbf y-\epsilon\|\mathbf y\|^2,\qquad\epsilon=\lambda_{\min}(\mathbf D_d)>0.\qquad\square
  \label{eq:os-impedance-sop}
\end{equation}
\end{proof}

弹簧的对称交叉项相消，与 \cref{ch:control} 中 $\dot{\mathbf M}-2\mathbf C$ 斜对称使科氏项消失（\cref{prop:ctrl-robot-props}）同源——能量观点贯穿机器人交互控制。无外力（$\mathbf f_{\rm ext}=\mathbf 0$）时 $\dot V=-\dot{\tilde{\mathbf x}}^\top\mathbf D_d\dot{\tilde{\mathbf x}}\le0$ 只半负定，由 LaSalle（\cref{ch:lyapunov-basics}）：$\dot V=0\Rightarrow\dot{\tilde{\mathbf x}}=\mathbf 0\Rightarrow\mathbf K_d^x\tilde{\mathbf x}=\mathbf 0\Rightarrow\tilde{\mathbf x}=\mathbf 0$，故渐近稳定。有外力时平衡点移到 $\tilde{\mathbf x}_{\rm eq}=(\mathbf K_d^x)^{-1}\mathbf f_{\rm ext}$（弹簧平衡位移）——这正是阻抗控制的设计意图。

\subsection{正实引理与 KYP：PR 与 SPR 的区分}
\label{sec:os-positive-real}

\begin{theorem}[正实引理 / KYP，频域无源判据]
\label{thm:os-kyp}
线性时不变系统 $H(s)=\mathbf C(s\mathbf I-\mathbf A)^{-1}\mathbf B+\mathbf D$ 无源，当且仅当 $H(s)$ \emph{正实}（PR）：在 $\mathrm{Re}(s)>0$ 上解析，且 $\mathrm{Re}[H(j\omega)]\ge0,\forall\omega$。KYP 引理给出其状态空间等价：存在 $\mathbf P=\mathbf P^\top\succ0,\mathbf L,\mathbf W$ 使 $\mathbf A^\top\mathbf P+\mathbf P\mathbf A=-\mathbf L^\top\mathbf L$、$\mathbf P\mathbf B=\mathbf C^\top-\mathbf L^\top\mathbf W$、$\mathbf W^\top\mathbf W=\mathbf D+\mathbf D^\top$。
\end{theorem}

对阻抗传递函数 \cref{eq:os-Zs}，
\begin{equation}
  \mathrm{Re}[Z(j\omega)]=\mathrm{Re}\Bigl[M_d j\omega+D_d-\tfrac{jK_d}{\omega}\Bigr]=D_d\ge0,
  \label{eq:os-ReZ}
\end{equation}
实部恒等于 $D_d$。须精确区分两个概念：$Z(s)$ 在 $s=0$ 有一个\emph{单极点}（来自弹簧项 $K_d/s$ 的积分器），故它是\emph{含原点单极点的正实}（PR）函数，而\emph{非}严格正实（SPR）。SPR 的标准定义要求对某 $\varepsilon>0$，$H(s-\varepsilon)$ 仍 PR（等价地在 $\mathrm{Re}(s)\ge0$ 上 $\mathrm{Re}[H]>0$），这要求函数在整条虚轴上解析——$Z$ 的原点极点恰好破坏这一点。从无源性看，$K_d/s$ 对应的积分器（弹簧）是\emph{无损}存储元件、$D_d>0$ 保证耗散，系统整体仍无源、但仅是 PR 不是 SPR。这一 PR/SPR 区分在鲁棒控制部的 IQC 统一判据里再次出现（\cref{ch:robust-hinf} 引用此处，不重证）。

\subsection{互联三定理与 Colgate--Hogan 耦合稳定}
\label{sec:os-interconnection}

\begin{theorem}[无源系统互联]
\label{thm:os-interconnect}
设 $\Sigma_1,\Sigma_2$ 无源、存储函数 $V_1,V_2$。则：
\begin{enumerate}
  \item \textbf{并联}（同输入、输出相加）仍无源；
  \item \textbf{级联}（$\Sigma_1$ 输出作 $\Sigma_2$ 输入）\emph{不一定}无源；
  \item \textbf{负反馈互联}：若 $\Sigma_1,\Sigma_2$ 均\emph{严格输出无源}，则反馈互联 $L_2$ 稳定；若两者再满足\emph{零状态可检测}（$\mathbf y_i\equiv\mathbf 0\Rightarrow$ 状态 $\to\mathbf 0$）且存储函数正定，则原点渐近稳定。\footnote{严格输出无源单独只给 $L_2$/有界性；渐近稳定还需零状态可检测——见 Khalil, \emph{Nonlinear Systems} 第 3 版定理 6.3。本章阻抗/环境互联的存储函数 \cref{eq:os-storage} 关于全状态正定且 $\mathbf y=\dot{\tilde{\mathbf x}}\equiv\mathbf 0$ 经动力学迫使 $\tilde{\mathbf x}\to\mathbf 0$，该前提自动满足，故下文沿用 LaSalle 直接得渐近稳定。}
\end{enumerate}
\end{theorem}
\begin{proof}
(1) 取 $V=V_1+V_2$，$\dot V=\dot V_1+\dot V_2\le\mathbf u^\top\mathbf y_1+\mathbf u^\top\mathbf y_2=\mathbf u^\top(\mathbf y_1+\mathbf y_2)$。
(2) 级联改变因果链，端口功率簿记断裂——$\Sigma_2$ 的输出不再与 $\Sigma_1$ 的输入直接挂钩，故无源性不传递（这正是导纳串联结构需逐层分析的根源）。
(3)（Colgate--Hogan 1988\cite{colgate1988robust}）设 $\dot V_i\le\mathbf u_i^\top\mathbf y_i-\epsilon_i\|\mathbf y_i\|^2$，负反馈条件 $\mathbf u_1=-\mathbf y_2,\mathbf u_2=\mathbf y_1$。取 $V=V_1+V_2$：
\[
  \dot V\le(-\mathbf y_2)^\top\mathbf y_1-\epsilon_1\|\mathbf y_1\|^2+\mathbf y_1^\top\mathbf y_2-\epsilon_2\|\mathbf y_2\|^2.
\]
关键一步：$(-\mathbf y_2)^\top\mathbf y_1+\mathbf y_1^\top\mathbf y_2=0$（标量转置等于自身），故
\begin{equation}
  \dot V\le-\epsilon_1\|\mathbf y_1\|^2-\epsilon_2\|\mathbf y_2\|^2\le0.
  \label{eq:os-feedback-stable}
\end{equation}
$\dot V$ 负半定，由 LaSalle 在 $\dot V=0$ 集上 $\mathbf y_1=\mathbf y_2=\mathbf 0\Rightarrow\mathbf u_1=\mathbf u_2=\mathbf 0$，系统回到原点，渐近稳定。
\end{proof}

落到力控：取 $\Sigma_1$ = 机器人 + 阻抗控制器（端口 $(\mathbf f,\mathbf v)$）、$\Sigma_2$ = 环境，耦合条件即 Newton 第三定律 $\mathbf f_{\rm robot}=-\mathbf f_{\rm env}$ 与速度连续 $\mathbf v_{\rm robot}=\mathbf v_{\rm env}$。于是只要阻抗控制器严格输出无源（$\mathbf D_d\succ0$），\emph{无论碰到什么被动环境——从软海绵到硬钢板——都保证稳定}，不需要知道环境参数。这是 Colgate--Hogan 的工程意义，也是把 \cref{thm:os-impedance-passive} 推到“耦合系统”的关键一步。

\begin{pitfall}[概念误区：把“无源”等同于“稳定”]
无源性是互联系统稳定的\emph{充分}条件，不是必要条件。非无源系统在特定环境下仍可能稳定（环境提供了足够阻尼）；反之，无源性保证的是\emph{对任意被动环境}的稳定——这是很强的鲁棒保证。追求无源性的理由正是它\emph{不依赖环境参数的精确知识}。
\end{pitfall}

\subsection{离散化破坏无源性}
\label{sec:os-discretization}

理论上无源的控制器，数字实现时未必无源：零阶保持引入相位滞后 $e^{-j\omega T/2}$（$T$ 为采样周期），等效成一个有源元件。Colgate 与 Schenkel（1997）给出采样数据系统无源的充分条件\cite{colgate1997passivity}：对“虚拟墙”阻抗，物理阻尼 $B$ 与虚拟阻尼 $D_d$ 须满足
\begin{equation}
  K_d<\frac{2(B-D_d)}{T},
  \label{eq:os-discrete-passive}
\end{equation}
即采样越慢、可用刚度上限越低。\textbf{注意虚拟阻尼 $D_d$ 以\emph{减号}进入}（等价条件 $B>K_dT/2+D_d$）：离散实现里速度常由差分估计、带一拍延迟，使虚拟阻尼在采样系统中\emph{反而注能}而非纯耗散，须由物理阻尼 $B$ 额外补偿——这是 Colgate--Schenkel 的反直觉核心（$B>K_dT/2-D_d$ 则是更弱的\emph{稳定}条件，非无源）。工程经验是采样率取力控带宽的 $10$ 倍以上。该结果常被误记为发表在 ASME 期刊，实际出处为 \emph{Journal of Robotic Systems}（Wiley）$14(1)$:$37$--$47$, 1997。若延迟来自通信/遥操作通道而非控制器本身，宜改用时域无源方法（TDPA）按离散端口功率累积能量、能量赤字时注入附加阻尼——TDPA 的工程实现属遥操作专题，本章只指针。

% ════════════════════════════════════════════════════════════════
\section{SE(3) 几何阻抗}
\label{sec:os-se3}

\paragraph{动机：姿态不是可以相减的向量。}
前几节的 $\tilde{\mathbf x}$ 暗含“误差能相减”。位置确可（$\mathbb R^3$ 是平坦的），但姿态活在弯曲的 $\SO(3)$ 上，简单相减既无几何意义又会在大角度时给出错误的恢复力矩。本节把阻抗几何正确地搬到 SE(3)，是本章记号最敏感的一幕——须严格沿用 \cref{ch:lie} 的右扰动与 $\se(3)=[\boldsymbol\rho;\boldsymbol\phi]$ 排序。

\subsection{欧拉角阻抗为何错}
\label{sec:os-euler-wrong}

用欧拉角（RPY）定义姿态误差有两个致命问题：\emph{万向锁}——pitch 接近 $\pm90^\circ$ 时表示退化，刚度矩阵在某些方向“消失”；\emph{非线性空间}——欧拉角空间不是线性空间，两旋转的“中点”不是欧拉角平均，直接对欧拉角做 PD 会在大角度偏转时产生非物理的恢复力矩，甚至让末端走“长路”（绕经万向锁奇异）而非“短路”。

\subsection{旋转误差的三种表示}
\label{sec:os-rot-error}

\begin{table}[H]\centering\small
\caption{旋转误差的三种表示与尺度对照（小角度极限）}
\label{tab:os-rot-error}
\begin{tabular}{p{0.24\textwidth} p{0.36\textwidth} p{0.30\textwidth}}
\toprule
表示 & 定义 & 小角度尺度 \\
\midrule
精确对数（本书主） & $\mathbf e_R=\mathrm{Log}(\mathbf R^\top\mathbf R_d)$（右扰动机体系） & $\theta\,\mathbf n$ \\
Caccavale 弦距 & $\mathbf e_R^{\rm cac}=\tfrac12\mathrm{vee}(\mathbf R_d^\top\mathbf R-\mathbf R^\top\mathbf R_d)$ & $\sin\theta\,\mathbf n$ \\
四元数虚部 & $\mathbf e_R^{\rm quat}=\mathrm{vec}(q^{-1}q_d)$（Hamilton） & $(\theta/2)\,\mathbf n$ \\
\bottomrule
\end{tabular}
\end{table}

本书\emph{精确式}用对数映射 $\mathbf e_R=\mathrm{Log}(\mathbf R^\top\mathbf R_d)$（\cref{ch:lie} 的右扰动，机体系旋转偏差），它是 GTSAM 等求解器里 $\SO(3)$ 残差的标准选择。Caccavale 等（1999）的弦距式 $\mathbf e_R^{\rm cac}=\tfrac12\mathrm{vee}(\mathbf R_d^\top\mathbf R-\mathbf R^\top\mathbf R_d)$\cite{caccavale1999sixdof} 是 $\mathrm{Log}$ 的一阶近似、计算更省，对应 Lyapunov 函数 $V_R=\tfrac12\mathrm{tr}(\mathbf I-\mathbf R_d^\top\mathbf R)$。三者方向一致但尺度不同（$\theta\,\mathbf n$ / $\sin\theta\,\mathbf n$ / $(\theta/2)\,\mathbf n$），故\emph{替换公式时不能直接沿用原旋转刚度}：近似相对误差为 $1-\sin\theta/\theta$，$14^\circ$ 时约 $1\%$、$30^\circ$ 时约 $4.5\%$；姿态误差可能超过 $15^\circ$ 时应改用精确对数（Rodrigues 逆 $\mathrm{Log}(\mathbf R_{\rm err})=\tfrac{\theta}{2\sin\theta}(\mathbf R_{\rm err}-\mathbf R_{\rm err}^\top)^\vee$，$\theta=\arccos\tfrac{\mathrm{tr}(\mathbf R_{\rm err})-1}{2}$）。

\subsection{四元数双覆盖}
\label{sec:os-quaternion}

四元数双覆盖（$q$ 与 $-q$ 表示同一旋转）必须处理：实现时若 $q\cdot q_d<0$ 则翻号 $q\leftarrow-q$，确保取\emph{最短路}。本书用 Hamilton 约定（标量在前），误差四元数取 $q^{-1}q_d$、误差向量取其虚部；实现库（libfranka/Eigen）标量在末，引用其 \texttt{coeffs()} 结果时须换序。最短路处理避免了“绕远路”的非物理恢复力矩，与 \cref{sec:os-euler-wrong} 的告诫同因。

\subsection{Stramigioli 势能：弦距 vs 测地}
\label{sec:os-stramigioli}

Stramigioli（2001）的几何 SE(3) 弹簧\cite{stramigioli2001modeling} 给出一个坐标无关的势能
\begin{equation}
  V_{\rm SE3}=\tfrac12 k_p\|\mathbf p-\mathbf p_d\|^2+k_r\,\mathrm{tr}(\mathbf I-\mathbf R_d^\top\mathbf R),
  \label{eq:os-se3-potential}
\end{equation}
其梯度给出几何正确的恢复力/力矩：$\mathbf F_{\rm spring}=-\nabla_{\mathbf p}V=-k_p(\mathbf p-\mathbf p_d)$（线性弹簧，与欧氏同）、$\boldsymbol\tau_{\rm spring}=-k_r\mathbf e_R$（旋转弹簧，用 Caccavale 误差）。须辨析：$\mathrm{tr}(\mathbf I-\mathbf R_d^\top\mathbf R)=2(1-\cos\theta)$ \emph{不是} $\SO(3)$ 上的\emph{测地}距离平方（后者为 $\theta^2=\|\mathrm{Log}(\mathbf R_d^\top\mathbf R)\|^2$），而是一种\emph{弦距}度量；二者小角度下一致（$2(1-\cos\theta)\approx\theta^2$），大角度有差异。弦距的好处是计算简单、梯度仍给出正确的恢复力矩方向，故工程上常用。

\paragraph{$6\times6$ 耦合刚度的量纲。}
按本书 $\se(3)=[\boldsymbol\rho;\boldsymbol\phi]$ 排序，刚度矩阵 $\mathbf K_d^x=\bigl[\begin{smallmatrix}\mathbf K_{\rho\rho}&\mathbf K_{\rho\phi}\\\mathbf K_{\phi\rho}&\mathbf K_{\phi\phi}\end{smallmatrix}\bigr]$ 四块量纲不同：$\mathbf K_{\rho\rho}$（N/m，位移生力）、$\mathbf K_{\phi\phi}$（Nm/rad，旋转生力矩）、$\mathbf K_{\rho\phi}$（N/rad，旋转生力，耦合）、$\mathbf K_{\phi\rho}$（Nm/m，位移生力矩，耦合）。多数实现取 $\mathbf K_{\rho\phi}=\mathbf K_{\phi\rho}^\top=\mathbf 0$（无耦合），避免量纲混合与数值条件化问题；斜面上的柔顺操作等任务确需耦合时，必须注意数值缩放。

\subsection{SE(3) 完整阻抗控制律}
\label{sec:os-se3-law}

把以上整合，定义广义位姿误差与速度误差（按本书排序，\emph{平移在前}）
\begin{equation}
  \mathbf e=\begin{bmatrix}\mathbf e_p\\\mathbf e_R\end{bmatrix}\in\mathbb R^6,\quad
  \mathbf e_p=\mathbf p-\mathbf p_d,\quad
  \mathbf e_R=\mathrm{Log}(\mathbf R^\top\mathbf R_d)\ (\text{或 Caccavale}),\qquad
  \dot{\mathbf e}=\mathbf J\dot{\mathbf q}-\begin{bmatrix}\dot{\mathbf p}_d\\\boldsymbol\omega_d\end{bmatrix},
  \label{eq:os-se3-error}
\end{equation}
其中 $\boldsymbol\omega$ 与 $\boldsymbol\omega_d$ 须在\emph{同一坐标系}表达。SE(3) 阻抗的完整控制律为
\begin{equation}
  \boxed{\;\boldsymbol\tau=\mathbf J^\top\bigl[\boldsymbol\Lambda(\ddot{\mathbf x}_d-\mathbf K_d^x\mathbf e-\mathbf D_d\dot{\mathbf e})+\boldsymbol\mu+\mathbf p\bigr]+\mathbf N^\top\bigl[-\mathbf K_{\rm null}(\mathbf q-\mathbf q_{\rm mid})-\mathbf D_{\rm null}\dot{\mathbf q}\bigr]\;}
  \label{eq:os-se3-impedance}
\end{equation}
六项各有物理含义：惯量前馈 $\boldsymbol\Lambda\ddot{\mathbf x}_d$（跟踪期望加速度）、弹簧力 $-\boldsymbol\Lambda\mathbf K_d^x\mathbf e$（位姿偏差恢复）、阻尼力 $-\boldsymbol\Lambda\mathbf D_d\dot{\mathbf e}$（抑制接触振荡）、科氏补偿 $\boldsymbol\mu$、重力补偿 $\mathbf p$、零空间正则化（关节回中/避奇异）。其数学结构与 Hogan 一行方程 $\mathbf M_d\ddot{\mathbf e}+\mathbf D_d\dot{\mathbf e}+\mathbf K_d^x\mathbf e=\mathbf f_{\rm ext}$ 完全一致，区别只在 $\mathbf e$ 含几何正确的旋转误差、$\boldsymbol\mu,\mathbf p$ 确保真机精确实现、$\mathbf N^\top$ 利用冗余。当取 $\mathbf M_d=\boldsymbol\Lambda$ 且 $\ddot{\mathbf x}_d=\mathbf 0$（稳态/慢变）时，律大幅简化为 $\boldsymbol\tau=\mathbf J^\top(-\mathbf K_d^x\mathbf e-\mathbf D_d\dot{\mathbf e}+\boldsymbol\mu+\mathbf p)+\mathbf N^\top\boldsymbol\tau_{\rm null}$，这正是真臂笛卡尔阻抗的核心（落地见 \cref{sec:ap-impedance}）。

\begin{pitfall}[坐标系一致性：一个负号导致正反馈失稳]
误差向量与雅可比\emph{必须在同一参考系}计算。工程上常用 \texttt{LOCAL\_WORLD\_ALIGNED} 帧（原点在末端、坐标轴与世界对齐——注意不是“位置在世界、姿态在末端”）；姿态误差从机体系旋到世界系时有一个负号 $-(\mathbf R\,q_{\rm err}.\mathrm{vec})$，使其与位置误差同取 current$-$desired，后续 $-\mathbf K\mathbf e$ 才产生指向期望的恢复力矩。若漏掉这个负号，末端会\emph{远离}期望姿态——这是一个\emph{正反馈}，立即失稳。排查法：让末端在自由空间偏离期望姿态 $10^\circ$，观察是否收敛；发散就检查误差符号。
\end{pitfall}

% ════════════════════════════════════════════════════════════════
\section{变阻抗与无源性保持}
\label{sec:os-variable}

\paragraph{动机：多阶段任务要改参数，但改参数会“充电”。}
固定阻抗在单一场景工作良好，但“接近 $\to$ 接触 $\to$ 操作 $\to$ 退出”的多阶段任务需要在线调参。然而 \cref{thm:os-impedance-passive} 的无源性证明依赖 $\mathbf K_d^x$ 常值——一旦时变，证明的关键步骤失效，系统可能凭空获得能量而失稳。本节给出三种把无源性\emph{找回来}的理论工具。

\subsection{变阻抗为何危险：能量注入}
\label{sec:os-energy-injection}

设末端偏差 $\tilde{\mathbf x}$ 固定、刚度从 $K_1$ 跳到 $K_2>K_1$，弹性储能从 $\tfrac12K_1\tilde x^2$ 跳到 $\tfrac12K_2\tilde x^2$，能量跳变 $\Delta V=\tfrac12(K_2-K_1)\tilde x^2>0$——这部分能量\emph{不来自外力、不来自电机，而是“换弹簧”凭空注入的}。严格地，对时变 $\mathbf K_d^x(t)$ 取存储函数 \cref{eq:os-storage}（$\mathbf K_d^x$ 换为 $\mathbf K_d^x(t)$），求导多出一项：
\begin{equation}
  \dot V=\dot{\tilde{\mathbf x}}^\top\mathbf f_{\rm ext}-\dot{\tilde{\mathbf x}}^\top\mathbf D_d\dot{\tilde{\mathbf x}}+\underbrace{\tfrac12\tilde{\mathbf x}^\top\dot{\mathbf K}_d^x\tilde{\mathbf x}}_{\text{注入项}}.
  \label{eq:os-injection}
\end{equation}
注入项的符号由 $\dot{\mathbf K}_d^x$ 决定：$\dot{\mathbf K}_d^x\succ0$（升刚度）注入能量、可能失稳；$\dot{\mathbf K}_d^x\prec0$（降刚度）提取能量、安全；$\dot{\mathbf K}_d^x=\mathbf 0$ 回到标准无源。能量注入有三个来源：刚度升高（最常见也最危险）、阻尼降低、参考轨迹跳变。一条工程上的正确性洞察：判断注入功率应用对称化 $\tfrac12\tilde{\mathbf x}^\top\mathrm{sym}(\dot{\mathbf K}_d^x)\tilde{\mathbf x}$（$\mathrm{sym}(\mathbf A)=\tfrac12(\mathbf A+\mathbf A^\top)$），而非简单取 $\dot{\mathbf K}_d^x$ 的迹。

\begin{pitfall}[思维陷阱：以为“变得慢就安全”]
缓慢变化只是\emph{降低}了瞬时注入率 $\tfrac12\tilde{\mathbf x}^\top\dot{\mathbf K}_d^x\tilde{\mathbf x}$，并不消除它；长时间缓慢升刚度仍会积累大量注入能量。是否安全取决于\emph{注入功率与阻尼耗散功率的竞争}——若耗散率 $\dot{\tilde{\mathbf x}}^\top\mathbf D_d\dot{\tilde{\mathbf x}}$ 始终大于注入率则安全。这正是下面 KB 条件量化的内容。另注意：降刚度虽提取能量，但若\emph{同时降阻尼}，耗散能力下降，极端情况下降刚度+降阻尼会进入欠阻尼振荡——降刚度时仍须检查阻尼比。
\end{pitfall}

\subsection{能量罐（Ferraguti--Secchi 2013）}
\label{sec:os-energytank}

\begin{theorem}[能量罐保无源]
\label{thm:os-energytank}
引入虚拟能量罐状态 $x_t\in\mathbb R$、罐能量 $T=\tfrac12 x_t^2\ge0$，扩展存储函数 $H=\tfrac12\dot{\tilde{\mathbf x}}^\top\mathbf M_d\dot{\tilde{\mathbf x}}+\tfrac12\tilde{\mathbf x}^\top\mathbf K_d^x(t)\tilde{\mathbf x}+T$。若设计罐动力学使
\begin{equation}
  \dot T=-\tfrac12\tilde{\mathbf x}^\top\dot{\mathbf K}_d^x\tilde{\mathbf x}+\beta\,\dot{\tilde{\mathbf x}}^\top\mathbf D_d\dot{\tilde{\mathbf x}},\qquad\beta\in[0,1),
  \label{eq:os-tank-dyn}
\end{equation}
则 $\dot H\le\dot{\tilde{\mathbf x}}^\top\mathbf f_{\rm ext}$，即以 $(\dot{\tilde{\mathbf x}},\mathbf f_{\rm ext})$ 为端口的扩展系统无源。
\end{theorem}
\begin{proof}
四步。\textbf{Step 1}：取上述 $H$。\textbf{Step 2}：求导并代入 $\mathbf M_d\ddot{\tilde{\mathbf x}}=\mathbf f_{\rm ext}-\mathbf D_d\dot{\tilde{\mathbf x}}-\mathbf K_d^x\tilde{\mathbf x}$、对称交叉项相消，得
\[
  \dot H=\dot{\tilde{\mathbf x}}^\top\mathbf f_{\rm ext}-\dot{\tilde{\mathbf x}}^\top\mathbf D_d\dot{\tilde{\mathbf x}}+\tfrac12\tilde{\mathbf x}^\top\dot{\mathbf K}_d^x\tilde{\mathbf x}+\dot T.
\]
\textbf{Step 3}：代入 \cref{eq:os-tank-dyn}，刚度注入项 $+\tfrac12\tilde{\mathbf x}^\top\dot{\mathbf K}_d^x\tilde{\mathbf x}$ 与 $\dot T$ 中的 $-\tfrac12\tilde{\mathbf x}^\top\dot{\mathbf K}_d^x\tilde{\mathbf x}$ \emph{完全抵消}，余 $\dot H=\dot{\tilde{\mathbf x}}^\top\mathbf f_{\rm ext}-(1-\beta)\dot{\tilde{\mathbf x}}^\top\mathbf D_d\dot{\tilde{\mathbf x}}$。\textbf{Step 4}：$\mathbf D_d\succ0$、$0\le\beta<1\Rightarrow\dot H\le\dot{\tilde{\mathbf x}}^\top\mathbf f_{\rm ext}$。
\end{proof}

数学上完美，工程上多一个约束 $T\ge0$：升刚度时罐能量减少，减到零就不能继续升。实现引入阀门函数 $\sigma,u\in\{0,1\}$——当需升刚度且本周期扣除后 $T-P_{\rm port}\Delta t\ge T_{\min}$ 时 $u=1$（从罐取能量）、否则 $u=0$ 并\emph{冻结} $\mathbf K_d^x$（$\dot{\mathbf K}_d^x:=\mathbf 0$，自动回到固定阻抗的无源情形）；$T<T_{\max}$ 时 $\sigma=1$ 允许充罐、罐满则停。

\begin{insight}
能量罐的精髓是\emph{把能量守恒从全局性质变成可管理的账本}。传统方法要求“系统总能量不增”——全局约束、难工程化；能量罐转化为“罐余额不为负”——局部状态约束，可用简单 if--else 实现。升刚度注入的能量不是凭空出现，而是从罐中\emph{转移}过来：总能量 $H$ 没增，只在子系统间重新分配。代价是引入保守性——罐空时冻结 $\mathbf K_d^x$ 暂时牺牲性能。这与 UFIC（\cref{sec:os-causality}）保证任意被动环境稳定用的是同一机制。
\end{insight}

\subsection{KB 阻尼下界（Kronander--Billard 2016）}
\label{sec:os-kb}

能量罐需在线维护罐状态；另一条更轻的路是 Kronander 与 Billard（2016）给出的\emph{离线}阻尼下界。

\begin{theorem}[KB 阻尼条件]
\label{thm:os-kb}
对零输入系统 $\mathbf M_d\ddot{\tilde{\mathbf x}}+\mathbf D_d(t)\dot{\tilde{\mathbf x}}+\mathbf K_d^x(t)\tilde{\mathbf x}=\mathbf 0$，若对所有 $t$
\begin{equation}
  \dot{\mathbf K}_d^x(t)+\alpha\dot{\mathbf D}_d(t)-2\alpha\mathbf K_d^x(t)\preceq\mathbf 0,\qquad
  \alpha\le\min_t\frac{\lambda_{\min}(\mathbf D_d(t))}{\lambda_{\max}(\mathbf M_d)},
  \label{eq:os-kb-condition}
\end{equation}
则系统指数稳定，$V(t)\le V(0)e^{-2\alpha t}$\cite{kronander2016stability}。
\end{theorem}
\begin{proof}[证明骨架]
取增广 Lyapunov 函数 $V_a=V+\alpha\dot{\tilde{\mathbf x}}^\top\tilde{\mathbf x}$（$V$ 为标准存储函数）。增广项引入位移—速度交叉耦合，使分析能同时利用二者信息。求 $\dot V_a$ 并代入 $\mathbf M_d\ddot{\tilde{\mathbf x}}=-\mathbf D_d\dot{\tilde{\mathbf x}}-\mathbf K_d^x\tilde{\mathbf x}$，整理出关于 $\dot{\tilde{\mathbf x}},\tilde{\mathbf x}$ 的二次型。要求 $\dot V_a+2\alpha V_a\le0$：速度二次型非正给出第一个条件（阻尼主导惯量 $\alpha\le\lambda_{\min}(\mathbf D_d)/\lambda_{\max}(\mathbf M_d)$）、位移二次型非正给出矩阵不等式 \cref{eq:os-kb-condition}。对足够小的 $\alpha$，$V_a$ 与 $V$ 等价正定，故 $V(t)\le c\,V(0)e^{-2\alpha t}$，指数稳定。
\end{proof}

当 $\mathbf D_d$ 常值（$\dot{\mathbf D}_d=\mathbf 0$），条件简化为 $\dot{\mathbf K}_d^x\preceq2\alpha\mathbf K_d^x$——\emph{允许}升刚度，只要升速率不超过 $2\alpha k_i$（与当前刚度成正比）。逐轴量纲自洽的工程形式为：第 $i$ 轴若升刚度，由 $\alpha_{{\rm req},i}=\max(0,\dot k_i)/(2k_{i,\min})$ 得所需最小衰减率（$\mathrm s^{-1}$），再由 $\alpha\le d_i/m_i$ 得阻尼下界，并与常规阻尼比要求取大者：
\begin{equation}
  d_{i,\min}=\max\Bigl(2\zeta_{\min}\sqrt{k_{i,\max}m_i},\ \frac{m_i\max(0,\dot k_i)}{2k_{i,\min}}\Bigr)\quad[\text{Ns/m}].
  \label{eq:os-kb-engineering}
\end{equation}
此处 $k_{i,\min}$ 取过渡区间内最小正刚度，避免在近零刚度时把允许增长率估得过大。须注意不能把 $\dot k/(2\alpha)$ 直接当阻尼——它的量纲是 N/m（刚度）而非 Ns/m（阻尼）。

\begin{table}[H]\centering\small
\caption{能量罐与 KB 阻尼条件的对照}
\label{tab:os-tank-vs-kb}
\begin{tabular}{p{0.22\textwidth} p{0.36\textwidth} p{0.34\textwidth}}
\toprule
特性 & 能量罐（Ferraguti） & KB 条件（Kronander--Billard） \\
\midrule
保证性质 & 端口无源（罐能量非负） & 零输入指数稳定 / 阻尼下界 \\
类型 & 在线状态监控 & 离线设计约束 \\
在线计算量 & 中（每周期更新罐） & 零（设计阶段完成） \\
保守性 & 低（允许更大变化） & 中（可能过度阻尼） \\
RL 兼容 & 好（可任意输出 $\mathbf K_d^x$） & 差（须同时输出满足条件的 $\mathbf D_d$） \\
\bottomrule
\end{tabular}
\end{table}

\emph{两层最佳实践}：用 KB 条件指导参数设计（确保大多数时间不需罐干预），用能量罐作运行时安全网（处理 KB 覆盖不到的极端情况）——两层保护优于单层。强化学习可输出变阻抗 $\mathbf K_d^x(t)$，但须配 KB 或能量罐保无源；学习型变阻抗的安全保证留给 \cref{ch:force-wbc} 与强化学习运控部。变阻抗安全也与规划部的控制屏障函数（CBF）安全证书互通：CBF 把“无源/能量预算”表达为安全约束，二者在“能量—安全证书”视角下统一。

% ════════════════════════════════════════════════════════════════
\section{接触动力学与力/速度椭球}
\label{sec:os-contact-ellipsoid}

\paragraph{动机：接触瞬间最危险，构型决定施力能力。}
前面把稳态接触讲透，但接触\emph{瞬间}的冲击与构型相关的施力能力同样是力控品质的根。本节给出接触过渡的冲击模型、超越线性弹簧的 Hunt--Crossley 模型、摩擦锥，以及力/速度椭球——并接上无传感力估计，把无源性的“安全”落到“感知”。

\subsection{接触过渡冲击}
\label{sec:os-impact}

牛顿碰撞模型给出接触前后法向速度的关系
\begin{equation}
  \dot x_n^+=-c_r\,\dot x_n^-,\qquad c_r\in[0,1],
  \label{eq:os-newton}
\end{equation}
$c_r$ 为恢复系数（完全弹性 $c_r=1$、完全塑性 $c_r=0$）。法向有效质量是构型依赖的：
\begin{equation}
  m_{\rm eff}=\frac{1}{\mathbf n^\top\boldsymbol\Lambda^{-1}\mathbf n},
  \label{eq:os-meff}
\end{equation}
$\mathbf n$ 为接触法向、$\boldsymbol\Lambda$ 是 \cref{sec:os-opspace-dynamics} 的操作空间惯量——同一构型在不同法向上“感觉重量”不同。须强调 $c_r$ 不确定性可达 $\pm30\%$（强依赖材料与撞击速度），故由 \cref{eq:os-impact} 估的冲击力数值只能作量级参考。工程上用四阶段平滑过渡（接近 $\to$ 接触 $\to$ 稳态 $\to$ 退出），例如用五次多项式平滑 $\mathbf K_d^x$ 的轨迹、令端点的一、二阶导为零，使 $\dot{\mathbf K}_d^x$ 在切换点连续（与 \cref{sec:os-variable} 的注入项控制一致）。

\subsection{Hunt--Crossley 接触模型}
\label{sec:os-hunt-crossley}

线性 Kelvin--Voigt 模型 $f_c=K\delta+D\dot\delta$ 有个非物理缺陷：分离前（$\dot\delta<0$）会预测\emph{受拉}力（接触面“吸住”）。Hunt 与 Crossley（1975）的非线性模型\cite{hunt1975coefficient} 修正了这一点：
\begin{equation}
  f_c=K_H\delta^n+D_H\delta^n\dot\delta,
  \label{eq:os-hunt-crossley}
\end{equation}
阻尼项含 $\delta^n$ 因子，故 $\delta\to0$（分离）时力平滑归零、不受拉。球—平面接触由 Hertz 理论给出 $n=3/2$；阻尼系数与恢复系数关联为 $D_H=\tfrac{3K_H(1-c_r)}{2\dot\delta^-}$（$\dot\delta^-$ 为撞击时穿透速率）。该模型把“恢复系数”这一离散碰撞概念解释为连续接触中的阻尼，是接触仿真与软接触力控的标准选择。

\subsection{摩擦锥}
\label{sec:os-friction}

切向接触受 Coulomb 摩擦约束，合法接触力须落在\emph{摩擦锥}内
\begin{equation}
  \|\mathbf f_t\|\le\mu f_n,
  \label{eq:os-friction-cone}
\end{equation}
$\mathbf f_t$ 为切向力、$f_n\ge0$ 为法向力、$\mu$ 为摩擦系数。在优化里常把二阶锥线性化为多边形金字塔，使约束成为线性、可纳入 QP。摩擦锥的接触约束优化（GIWC/LCP）见 \cref{ch:contact-mechanics}，足式 WBC 里同一摩擦锥的处理见 \cref{ch:quad_wbc}——同一数学、不同场景，本章只点出力控所需。

\subsection{力/速度/惯量椭球}
\label{sec:os-ellipsoid}

末端的施力与运动能力由雅可比的奇异值结构刻画。设 $\mathbf J$ 的奇异值 $\sigma_1\ge\cdots\ge\sigma_m$。\emph{速度椭球}是 $\|\dot{\mathbf q}\|\le1$ 时末端速度集合、\emph{力椭球}是 $\|\boldsymbol\tau\|\le1$ 时末端力集合：
\begin{equation}
  \mathcal E_v=\{\mathbf v:\mathbf v^\top(\mathbf J\mathbf J^\top)^{-1}\mathbf v\le1\},\qquad
  \mathcal E_f=\{\mathbf f:\mathbf f^\top\mathbf J\mathbf J^\top\mathbf f\le1\},
  \label{eq:os-ellipsoid}
\end{equation}
速度椭球半轴长 $=\sigma_i$、力椭球半轴长 $=1/\sigma_i$——二者半轴\emph{互为倒数}（对偶）。直觉如自行车齿轮比：某方向运动灵活（$\sigma_i$ 大）则施力困难（$1/\sigma_i$ 小），反之亦然。打磨任务“法向精密力 + 切向快速运动”遂成构型空间里的权衡。Yoshikawa（1985）操作度 $w=\sqrt{\det(\mathbf J\mathbf J^\top)}=\prod_i\sigma_i$ 衡量椭球体积（$w=0$ 即奇异）\cite{yoshikawa1985manipulability}；Khatib（1995）\emph{动态}操作度 $w_{\rm dyn}=\sqrt{\det(\mathbf J\mathbf M^{-1}\mathbf J^\top)}$ 把惯量分布也纳入\cite{khatib1995inertial}。操作空间惯量的条件数 $\kappa(\boldsymbol\Lambda)$ 直接影响力控品质——这正是 \cref{sec:os-dls} 要保证 $\sigma_{\min}$ 不过小的根据。

\begin{figure}[ht]
\centering
\begin{tikzpicture}[>=Stealth, font=\small, scale=1.0]
  % velocity ellipse: long horizontal, short vertical
  \draw[third,very thick] (-2.3,0) ellipse [x radius=1.7, y radius=0.6];
  \node[third] at (-2.3,1.0) {速度椭球 $\mathcal E_v$};
  \draw[->,third] (-2.3,0) -- (-0.6,0) node[midway,below]{$\sigma_1$};
  \draw[->,third] (-2.3,0) -- (-2.3,0.6) node[midway,left]{$\sigma_2$};
  % force ellipse: short horizontal, long vertical (dual)
  \draw[second,very thick] (2.3,0) ellipse [x radius=0.6, y radius=1.7];
  \node[second] at (2.3,2.05) {力椭球 $\mathcal E_f$};
  \draw[->,second] (2.3,0) -- (2.9,0) node[midway,below]{$1/\sigma_1$};
  \draw[->,second] (2.3,0) -- (2.3,1.7) node[midway,right]{$1/\sigma_2$};
  % duality arrow
  \draw[<->,thick,gray] (-0.45,-0.05) -- (1.65,-0.05);
  \node[gray] at (0.6,-0.4) {对偶（半轴互倒）};
\end{tikzpicture}
\caption{力/速度椭球的对偶：速度椭球长轴方向（易运动）正是力椭球短轴方向（难施力）。半轴长互为倒数 $\sigma_i\leftrightarrow1/\sigma_i$。}
\label{fig:os-ellipsoid}
\end{figure}

\subsection{无传感力估计：动量观测器}
\label{sec:os-observer}

许多机器人没有腕部力/力矩传感器；动量观测器（De Luca--Mattone 2003）用\emph{广义动量}估外力，无需测量噪声大的关节加速度。

\begin{theorem}[广义动量观测器]
\label{thm:os-observer}
对 $\mathbf M\ddot{\mathbf q}+\mathbf C\dot{\mathbf q}+\mathbf g=\boldsymbol\tau+\boldsymbol\tau_{\rm ext}$，定义广义动量 $\mathbf p=\mathbf M\dot{\mathbf q}$ 与残差
\begin{equation}
  \mathbf r=\mathbf K_I\Bigl[\mathbf p-\int_0^t(\boldsymbol\tau+\mathbf C^\top\dot{\mathbf q}-\mathbf g+\mathbf r)\,\mathrm d\tau'-\mathbf p(0)\Bigr],
  \label{eq:os-residual}
\end{equation}
则 $\mathbf r$ 满足一阶线性滤波 $\dot{\mathbf r}=\mathbf K_I(\boldsymbol\tau_{\rm ext}-\mathbf r)$，以时间常数 $\mathbf K_I^{-1}$ 指数收敛到外力矩 $\boldsymbol\tau_{\rm ext}$。
\end{theorem}
\begin{proof}[证明骨架]
求 $\dot{\mathbf p}=\mathbf M\ddot{\mathbf q}+\dot{\mathbf M}\dot{\mathbf q}$，代入动力学得 $\dot{\mathbf p}=\boldsymbol\tau+\boldsymbol\tau_{\rm ext}-\mathbf C\dot{\mathbf q}-\mathbf g+\dot{\mathbf M}\dot{\mathbf q}$。用斜对称性质 $\dot{\mathbf M}=\mathbf C+\mathbf C^\top$（即 $\dot{\mathbf M}-2\mathbf C$ 斜对称，\cref{prop:ctrl-robot-props}）化简得 $\dot{\mathbf p}=\boldsymbol\tau+\boldsymbol\tau_{\rm ext}+\mathbf C^\top\dot{\mathbf q}-\mathbf g$。对 \cref{eq:os-residual} 求导并代入 $\dot{\mathbf p}$，$\boldsymbol\tau,\mathbf C^\top\dot{\mathbf q},\mathbf g$ 项相消，得 $\dot{\mathbf r}=\mathbf K_I(\boldsymbol\tau_{\rm ext}-\mathbf r)$。
\end{proof}

\begin{insight}
精妙之处在于\emph{不需测 $\ddot{\mathbf q}$}：利用 $\mathbf p=\mathbf M\dot{\mathbf q}$ 与 $\dot{\mathbf M}=\mathbf C+\mathbf C^\top$，只需可直接测量的 $\mathbf q,\dot{\mathbf q},\boldsymbol\tau$ 就能估外力，比从 $\boldsymbol\tau_{\rm ext}=\mathbf M\ddot{\mathbf q}+\mathbf C\dot{\mathbf q}+\mathbf g-\boldsymbol\tau$（需 $\ddot{\mathbf q}$）稳健得多。
\end{insight}

把关节残差映射到任务空间 wrench，可用动力学一致伪逆 $\hat{\mathbf f}_{\rm ext}=\bar{\mathbf J}^\top\mathbf r=\boldsymbol\Lambda\mathbf J\mathbf M^{-1}\mathbf r$（注意 $\mathbf M^{-1}$ 已含在 $\bar{\mathbf J}^\top=\boldsymbol\Lambda\mathbf J\mathbf M^{-1}$ 中，勿重复乘），或 Moore--Penrose $\hat{\mathbf f}_{\rm ext}=(\mathbf J\mathbf J^\top)^{-1}\mathbf J\mathbf r$、近奇异时加 DLS。残差还可用于无 F/T 传感的标定（重力补偿三姿态最小二乘，原理见 \cref{ch:nlopt}）。基于残差的碰撞检测/隔离/识别/反应四阶段框架见 Haddadin 等（2017，\emph{IEEE Transactions on Robotics} $33(6)$:$1292$--$1312$，非 RAM）\cite{haddadin2017collision} 与 De Luca 等（2006）\cite{deluca2006collision}；ISO/TS 15066:2016 给出协作机器人功率—力限值（如前额准静态约 $130\,$N、胸部约 $140\,$N，瞬态约为准静态 $2$ 倍）作为安全约束的语义。\emph{完整}的碰撞反应策略与安全认证留给 \cref{ch:force-wbc} 与机械臂部实战。

% ════════════════════════════════════════════════════════════════
\section{桥接与小结}
\label{sec:os-bridge}

\paragraph{与他章的关系。}
本章是 \cref{ch:control} “机器人控制与交互”一节的严格化深化：那里给入门骨架（操作空间结论式、阻抗哲学、PD 无源），本章给完整推导（Khatib 投影、$\bar{\mathbf J}$/$\mathbf N$、频域 $Z(s)$、SE(3) 几何、无源性五定理、变阻抗保稳）。一般无源性/LaSalle/ISS 系统理论见 \cref{ch:lyapunov-stability}（本章只取接触交互所需）；正实/KYP 引理在本章证一次，鲁棒控制部 IQC 视角（\cref{ch:robust-hinf}）引用。向下，力位混合（Raibert--Craig 选择矩阵的完整版）、笛卡尔阻抗工程实现、WBC 多优先级 HQP/TSID 见 \cref{ch:force-wbc}；复合系统统一动力学与多模态 MPC 见 \cref{ch:unified-multimodal-mpc}；接触约束优化（摩擦锥/GIWC/LCP）见 \cref{ch:contact-mechanics}；操作空间惯量的空间向量算法见 \cref{ch:spatial-dynamics}。真臂落地（笛卡尔阻抗、$\mathbf J^+$ vs $\bar{\mathbf J}$ 的病态实测）见 \cref{sec:ap-impedance,sec:ap-illcond}；足式 WBC 的摩擦锥与多优先级见 \cref{ch:quad_wbc}。学习型变阻抗与遥操作二端口分别在强化学习运控部与机械臂部展开。本章讲透的“操作空间动力学 + 阻抗/导纳 + 无源性”，正是机械臂部理论章可直接 \cref{ch:operational-space} 复用、不必重建的地基。

\paragraph{与最优性的一句对照。}
无源性给的是“对任意被动环境鲁棒稳定”，最优控制给的是“某代价泛函最优”——二者视角互补：PD+重力补偿/阻抗靠无源保稳而不追最优，LQR（\cref{ch:lqr-lqg-riccati}）/HJB（\cref{ch:dp-hjb}）追最优而需模型。接触交互场景因环境未知，无源性的鲁棒保证常比最优性更切题。

\begin{table}[H]\centering\small
\caption{本章速查：操作空间公式、无源性判据、阻抗参数约束}
\label{tab:os-cheatsheet}
\begin{tabular}{p{0.30\textwidth} p{0.62\textwidth}}
\toprule
条目 & 关键式 / 判据 \\
\midrule
操作空间惯量 & $\boldsymbol\Lambda=(\mathbf J\mathbf M^{-1}\mathbf J^\top)^{-1}$ \\
动力学一致伪逆 & $\bar{\mathbf J}=\mathbf M^{-1}\mathbf J^\top\boldsymbol\Lambda$；$\bar{\mathbf J}^\top\mathbf M\bar{\mathbf J}=\boldsymbol\Lambda$、$\mathbf J\bar{\mathbf J}=\mathbf I$ \\
零空间投影 & $\mathbf N=\mathbf I-\bar{\mathbf J}\mathbf J$；$\mathbf N^2=\mathbf N$、$\mathbf J\mathbf N=\mathbf 0$、$\mathbf N^\top\mathbf M=\mathbf M\mathbf N$ \\
阻抗传递函数 & $Z(s)=M_d s+D_d+K_d/s$；$\mathrm{Re}[Z(j\omega)]=D_d\ge0$（PR，含原点单极点，非 SPR） \\
严格输出无源 & $\dot V\le\mathbf u^\top\mathbf y-\epsilon\|\mathbf y\|^2$，$\epsilon=\lambda_{\min}(\mathbf D_d)$ \\
Colgate--Hogan & 两端严格输出无源 $\Rightarrow$ 反馈互联渐近稳定（无模型） \\
临界阻尼 & $\mathbf D_d=2\mathbf M_d^{1/2}(\mathbf M_d^{-1/2}\mathbf K_d^x\mathbf M_d^{-1/2})^{1/2}\mathbf M_d^{1/2}$ \\
环境耦合 & $\omega_{\rm cl}=\sqrt{(K_d+K_e)/M_d}$、$f_{ss}=K_eK_d x_d/(K_d+K_e)$ \\
变阻抗注入 & $\tfrac12\tilde{\mathbf x}^\top\dot{\mathbf K}_d^x\tilde{\mathbf x}$；升刚度注入、降刚度提取 \\
KB 阻尼下界 & $\dot{\mathbf K}_d^x+\alpha\dot{\mathbf D}_d-2\alpha\mathbf K_d^x\preceq\mathbf 0$，$\alpha\le\lambda_{\min}(\mathbf D_d)/\lambda_{\max}(\mathbf M_d)$ \\
SE(3) 误差三表示 & $\mathrm{Log}(\mathbf R^\top\mathbf R_d)$（$\theta\mathbf n$）/ Caccavale（$\sin\theta\mathbf n$）/ 四元数虚部（$(\theta/2)\mathbf n$） \\
柔性关节刚度上界 & $K_{d,\rm eq}=K_\theta K_s/(K_\theta+K_s)\le K_s$（Ott 2008\cite{ott2008cartesian}） \\
\bottomrule
\end{tabular}
\end{table}

\subsection*{常见误解汇总}
\begin{itemize}
  \item \textbf{“无源 = 稳定”}：无源性是互联稳定的充分非必要条件；其价值在于对\emph{任意}被动环境的鲁棒保证（\cref{thm:os-interconnect}）。
  \item \textbf{“$\bar{\mathbf J}$ 与 $\mathbf J^+$ 只是精度差别”}：是物理正确性差别——$\mathbf J^+$ 让零空间力矩经惯量泄漏到末端，是“阻抗抖动”的根因（\cref{sec:os-jbar-vs-pinv}）。
  \item \textbf{“变阻抗变得慢就安全”}：缓慢只降注入率不消除它，长时间仍积累；安全取决于注入与耗散的竞争（\cref{sec:os-kb}）。
  \item \textbf{“$\mathrm{tr}(\mathbf I-\mathbf R_d^\top\mathbf R)$ 是测地距离平方”}：它是\emph{弦距}（$=2(1-\cos\theta)\neq\theta^2$），小角度才近似一致（\cref{sec:os-stramigioli}）。
  \item \textbf{“欧拉角做姿态阻抗也行”}：万向锁 + 非线性空间会给非物理恢复力矩、走长路（\cref{sec:os-euler-wrong}）。
\end{itemize}

\subsection*{延伸阅读}
操作空间动力学原始文献为 Khatib（1987）\cite{khatib1987operational}，惯量层框架与动态操作度见 Khatib（1995）\cite{khatib1995inertial}；阻抗控制三部曲为 Hogan（1985）\cite{hogan1985impedance}，主动刚度先驱 Salisbury（1980）\cite{salisbury1980active}、被动柔顺 RCC 见 Whitney（1982）\cite{whitney1982quasistatic}；耦合稳定与无源性为 Colgate--Hogan（1988）\cite{colgate1988robust}、采样无源 Colgate--Schenkel（1997）\cite{colgate1997passivity}；统一无源框架与双对角化见 Albu-Sch\"affer 等（2007）\cite{albuschaffer2007unified}、柔性关节系统化处理见 Ott（2008）\cite{ott2008cartesian}；变阻抗保稳的能量罐见 Ferraguti 等（2013）\cite{ferraguti2013tank}、阻尼下界见 Kronander--Billard（2016）\cite{kronander2016stability}、力—阻抗统一见 Haddadin--Shahriari（2024）\cite{shahriari2024unified}；几何 SE(3) 弹簧见 Stramigioli（2001）\cite{stramigioli2001modeling}、角/轴姿态阻抗见 Caccavale 等（1999）\cite{caccavale1999sixdof}；接触模型见 Hunt--Crossley（1975）\cite{hunt1975coefficient}、碰撞检测综述见 Haddadin 等（2017）\cite{haddadin2017collision}、动量观测器原理见 De Luca--Mattone（2003）\cite{deluca2003actuator}；操作空间/阻抗/力控的标准教材见 Spong 等\cite{spong2006robot}。
